沿一维实验室正 x 方向运动的探测器速度为 u≥0，质量为 m>0；质量 M>m 的行星沿同方向速度为 V>u。行星从后方追上探测器，发生理想一维完全弹性的折返散射，不发生碰撞破坏，初末相对距离都足够远以忽略引力势能。初步将 M 视为远大于 m。
(1) 在行星惯性系求散射前后探测器相对速度和实验室末速度；(2) 求探测器的动能增加，说明为何增益符号与 u、V 的相对大小有关；(3) 对有限 M 和 m 从动量、能量守恒求末速度；(4) 核验有限质量解的相对速度反号以及 M→∞ 极限，解释行星的能量损失。
